% ID: SEL-D0084-OA-E03
% Tanda S084.
Sean $B$ blanca, $N$ negra y $M$ marcada. Organizamos los recuentos:
\[
 \begin{array}{c|rr|r}
 &M^c&M&\text{Total}\\\hline
 B&25&75&100\\
 N&125&175&300\\\hline
 \text{Total}&150&250&400
 \end{array}
\]
\textbf{a)} $P(B)=100/400=\boxed{1/4}$.

\textbf{b)} Dentro del grupo de 250 bolas marcadas hay 75 blancas. Por tanto, $P(B\mid M)=75/250=\boxed{3/10}$.

\textbf{c)} La probabilidad conjunta de negra y marcada es $P(N\cap M)=175/400=\boxed{7/16}$.

\textbf{d)} Se cumple $P(B\cap M)=75/400=3/16$, mientras que $P(B)P(M)=(1/4)(250/400)=5/32$. Como son distintos, los sucesos son dependientes.
